Fix the score vector.  We first count the possible index sets without using any probabilistic premise.  Write a point of \(\mathbb R^2\) as \(v=(v_1,v_2)\).  Membership of \(X_i\) in the closed disk with center \((a,b)\) and radius \(r\) is equivalent to
\[
 2aX_{i1}+2bX_{i2}+s\ge \lVert X_i\rVert_2^2,
 \qquad s=r^2-a^2-b^2. \tag{1}
\]
Thus the traces on \(X_1,\ldots,X_n\) of Euclidean disks form a subcollection of the sign patterns generated by \(n\) affine hyperplanes in the three-dimensional parameter \((a,b,s)\).  A pattern defined by weak inequalities in (1) is also realized in an open cell: increase \(s\) by a sufficiently small positive amount, smaller than the minimum strictly negative margin among the finitely many excluded observations.  If there is no excluded observation, any positive increase works.

Let \(R_d(m)\) be the maximum number of open regions cut out by \(m\) affine hyperplanes in \(\mathbb R^d\).  Adding the last hyperplane leaves every old region intact or splits it into two, and a split can occur only in a region of the arrangement induced on that last hyperplane.  Consequently
\[
 R_d(m)\le R_d(m-1)+R_{d-1}(m-1),
 \qquad R_d(0)=R_0(m)=1. \tag{2}
\]
Induction on \(d+m\), using Pascal's identity, gives
\[
 R_d(m)\le\sum_{j=0}^d\binom mj. \tag{3}
\]
Equations (1)--(3) show that all disks induce at most \(\sum_{j=0}^3\binom nj\) subsets.  For fixed \(t\), intersection with the fixed set \(\{i:X_i\in\mathcal A_t\}\) cannot increase this number.  Taking the union over the two sides and over \(h\in\mathcal H_n\) proves
\[
 |\mathscr S_n|\le 2|\mathcal H_n|\sum_{j=0}^3\binom nj. \tag{4}
\]

We next prove the probabilistic assertions from the exact Gaussian model.  By \(\mathrm{ass:iid\text{-}sampling}\), the observations are independent.  By \(\mathcal A_0=\mathbb R^2\setminus\mathcal A_1\), the side containing \(X_i\) is \(T_i\).  Assumption \(\mathrm{ass:gaussian\text{-}errors}\) says that, for almost every score value on that side,
\[
 \mathcal L_P\{Y_i-g_{T_i,P}(X_i)\mid X_i\}=\mathcal N(0,\sigma^2). \tag{5}
\]
Independence of the sampled units implies conditional independence given \((X_1,\ldots,X_n)\).  Hence, for every deterministic \(S\in\mathscr S_n\) and \(\lambda\in\mathbb R\), multiplication of the exact conditional Gaussian moment-generating functions gives
\[
\begin{aligned}
 \mathbb E_P\left[\exp\left\{\lambda\sum_{i\in S}\varepsilon_i\right\}
       \middle|X_1,\ldots,X_n\right]
 &=\prod_{i\in S}\mathbb E_P\left[e^{\lambda\varepsilon_i}\mid X_i\right]\\
 &=\prod_{i\in S}e^{\sigma^2\lambda^2/2}
 =e^{|S|\sigma^2\lambda^2/2}. \tag{6}
\end{aligned}
\]
For \(u>0\), conditional Markov inequality applied to (6), followed by minimization at \(\lambda=u/(|S|\sigma^2)\), yields
\[
 \Pr_P\left\{\sum_{i\in S}\varepsilon_i>u\mid X_1,\ldots,X_n\right\}
 \le \exp\left\{-\frac{u^2}{2|S|\sigma^2}\right\}. \tag{7}
\]
Applying (7) to the negative sum and setting \(u=z\sigma\sqrt{|S|}\) gives
\[
 \Pr_P\left\{\frac{|\sum_{i\in S}\varepsilon_i|}{\sigma\sqrt{|S|}}>z
       \middle|X_1,\ldots,X_n\right\}\le2e^{-z^2/2}. \tag{8}
\]
The set \(\mathscr S_n\) is fixed after conditioning on the scores.  Summing (8) over it proves the asserted maximum tail.

Finally put \(m_S=|\mathscr S_n|\).  If \(m_S=0\), the stated convention proves the expectation claim.  If \(m_S\ge1\), put \(a=\sqrt{2\log(2m_S)}\).  Integrating the tail bound and using \(\int_a^\infty e^{-z^2/2}\,dz\le a^{-1}e^{-a^2/2}\) gives
\[
\begin{aligned}
 \mathbb E_P\left[\max_{S\in\mathscr S_n}
 \frac{|\sum_{i\in S}\varepsilon_i|}{\sigma\sqrt{|S|}}
 \middle|X_1,\ldots,X_n\right]
 &\le a+2m_S\int_a^\infty e^{-z^2/2}\,dz\\
 &\le a+a^{-1}\le a+1, \tag{9}
\end{aligned}
\]
because \(a\ge\sqrt{2\log2}>1\).  This is the claimed bound and completes the proof.